\section*{ВВЕДЕНИЕ}
\addcontentsline{toc}{section}{ВВЕДЕНИЕ}

Компилятор языка программирования Common Lisp представляет собой программную систему, осуществляющую трансляцию исходного кода, написанного на языке высокого уровня, в машинный код или байт-код виртуальной машины. В отличие от интерпретаторов, компилятор выполняет статический анализ и преобразование программы перед её исполнением, что открывает возможности для применения глубоких оптимизаций, повышающих эффективность выполнения.

Разработка собственного компилятора с поддержкой современных методов оптимизации позволяет не только достичь высокой производительности генерируемого кода, но и обеспечить полный контроль над процессом трансляции, адаптацию под конкретную среду исполнения, а также возможность исследования и внедрения эффективных приёмов реализации функциональных конструкций, таких как хвостовая рекурсия, замыкания и макросы.

Разработанный компилятор генерирует на основе исходного кода Common Lisp байт-код для выполнения на виртуальной машине. Поддерживаются следующие оптимизации: свёртка констант, удаление неиспользуемого кода, оптимизация хвостовой рекурсии, упрощение тривиальных условий, оптимизация арифметических выражений, встраивание функций.

\emph{Цель настоящей работы} -- реализация программной системы оптимизирующего компилятора Common Lisp. Для достижения поставленной цели необходимо решить \emph{следующие задачи:}
\begin{itemize}
\item изучение процесса компиляции языков программирования и функциональных языков в частности;
\item изучение методов оптимизации компиляторов;
\item проектирование компилятора с оптимизациями;
\item реализация компиляции S-выражения в промежуточную форму;
\item реализация генерации кода ассемблера;
\item реализация генерации байт-кода;
\item реализация оптимизации промежуточной формы.
\end{itemize}

\emph{Структура и объем работы.} Отчет состоит из введения, 4 разделов основной части, заключения, списка использованных источников, 2 приложений. Текст выпускной квалификационной работы равен \formbytotal{lastpage}{страниц}{е}{ам}{ам}.

\emph{Во введении} сформулирована цель работы, поставлены задачи разработки, описана структура работы, приведено краткое содержание каждого из разделов.

\emph{В первом разделе} на стадии описания технической характеристики предметной области приводится сбор информации о методах разработки компиляторов.

\emph{Во втором разделе} на стадии технического задания приводятся требования к разрабатываемому компилятору.

\emph{В третьем разделе} на стадии технического проектирования представлены проектные решения для компилятора.

\emph{В четвертом разделе} приводится список модулей и их функций, использованных при разработке компилятора, производится тестирование разработанной программной системы компилятора.

В заключении излагаются основные результаты работы, полученные в ходе разработки.

В приложении А представлен графический материал.
В приложении Б представлены фрагменты исходного кода. 
